% ============================================================
%  R. Nielsen, "Om S. Kierkegaards „mentale Tilstand“".
%  Nordisk Universitets-Tidsskrift, Fjerde Aargang, Tredie Hefte.
%  Kjøbenhavn 1858, pp. 1–29.
%
%  TRANSCRIPTION (Danish) — from the Google Books scan.
%  Scan: ~/bibliotek/Nielsen, Rasmus/Om_S_Kierkegaards_mentale_tilstand.pdf
%        (46 PDF pp., Google Books, Harvard College Library copy, shelf-mark
%        pencilled "Scan 6664.585.5"; 600 ppi bitonal. An identical earlier
%        download sits beside it as Kierkegaards_mentale_tilstand.pdf — same
%        46 pages, same image objects; either may be used.)
%
%  PAGE MAP (verified 2026-09-12: every body numeral read off the PDF text
%  layer in sequence, endpoints eye-confirmed on the images):
%        printed  1–29  =  PDF + 12   (PDF 13–41), uniform, no gaps, no dupes
%        PDF 1–10   Google/library front matter, the volume's half-title and
%                   title leaves.
%        PDF 11–12  the journal's FORORD and the lists of contributors and
%                   commissioners — NOT part of this article.
%        PDF 42     blank.  PDF 43–46  library date slip and back leaves.
%        Folio AND running head are both suppressed on printed p. 1.
%        Trap: the bottom right of pp. 1, 3, 17, 19 carries a SHEET SIGNATURE
%        (1, 1*, 2, 2*), not a folio. Folios proper sit at the OUTER TOP.
%
%  TYPEFACE: ANTIQUA throughout — a Didone roman — body, heads and footnotes
%        alike. There is NO ITALIC anywhere in the article, and no Fraktur and
%        no Greek. German quoted in the body is set in the same roman (p. 25,
%        "Ja, ich bin der Atheist und Gottlose!"). OCR pipeline accordingly
%        uses tesseract -l dan, with the PDF's own Google text layer as the
%        second witness (~0.80 type-level agreement, measured on pp. 5, 13, 23).
%
%  CONVENTIONS: diplomatic transcription; 1858 orthography as printed;
%  printer's errors transcribed as printed and logged in % comments at the
%  site, never silently corrected.
%
%   ö / ø  — BOTH SORTS OCCUR AND THE COMPOSITOR MIXES THEM. ö (o-diaeresis)
%        is the default and by far the commoner (hörte, sögte, förste, gjöre,
%        Öine, Berömmelse); but ø turns up, twice within twelve lines of an ö
%        on p. 19 ("under Øine", "Følgeblade") and again on p. 23 ("Øiet").
%        No OCR engine preserves the distinction — every one flattens it — so
%        each o-with-a-diacritic is an image decision. Transcribed as printed.
%   aa    — always aa, never å. Not one ring in the article (saaledes, maatte,
%        ogsaa, Aandsanstrængelser). æ is a true ligature throughout.
%   EMPHASIS is LETTERSPACING (Sperrsatz) ONLY, rendered \emph{}. Since the
%        article has no italic at all, every \emph{} in this file is a
%        letterspacing call. spacing.py finds them mechanically; it is an aid,
%        and the image is the authority.
%   QUOTES „…“ as printed: opening low, closing raised, UNEQUAL heights, never
%        ”. No guillemets anywhere in the 1858 printing (the modern spelling of
%        the title with » « is not the book's). NESTED QUOTATIONS REUSE THE
%        SAME PAIR (pp. 21, 23), so the quote-balance count in check.py is
%        unreliable on this text and a non-zero balance is not by itself a
%        defect — see the log below.
%   SPACED PUNCTUATION. The scan shows a thin space before , . : ; ? ! in
%        stretched lines ("Dage , hörte") and none in tight ones. It is
%        justification, not copy, and is NOT reproduced; the few places where
%        it survives in an unstretched short line are logged rather than set.
%   ELLIPSES are spaced full stops, four or five of them, and the count varies
%        (p. 1 opens with five). The printed count is reproduced.
%   NOT TRANSCRIBED: running heads, folios, sheet signatures, catchwords, and
%        the library's pencil marks and stamps. The running head's final
%        punctuation varies from page to page (p. 17 »Tilstand".«, p. 21
%        »Tilstand"«, p. 25 »Tilstand."«); recorded here as a datum, not set.
%
%  FOOTNOTES: two only, on printed pp. 17 and 25, both marked *), both set
%        smaller under a short left-hand rule, numbering restarting per page.
%        Hence \renewcommand{\thefootnote}{*)} rather than \fnsymbol (which is
%        math mode and gives U+2217) or footmisc[perpage] (which resets on the
%        TYPESET page, not the printed one).
%
%  STRUCTURE: none — no sections, no numbering, no ornaments. Rules occur only
%        on p. 1 (two, around the epigraph), pp. 17 and 25 (footnote rules) and
%        p. 29 (a short terminal rule). The long Kierkegaard quotations are NOT
%        block quotations: full measure, same size, ordinary paragraph indent,
%        delimited by „…“ and occasionally letterspaced for prominence (p. 2).
%        The article ends cleanly on p. 29 ("… i det förste Afsnit.") with no
%        signature, no dateline and no "Fortsættes".
%
%  STATUS: COMPLETE. pp. 1–29 transcribed and image-verified page by page,
%        2026-09-12, in two batches (1–15, 16–29). Every page was rendered and
%        read; doubtful glyphs were cropped at 900–1200 ppi.
%
%  ö / ø CENSUS (the distinction no OCR preserves; settled on the images).
%        FOUR ø in the whole article, all in the second half; every other
%        o-diacritic — some 300 of them — is ö:
%          p. 19  (Følgeblade   and   under Øine,   (the latter on the same
%                 line as höieste, which is the calibration)
%          p. 23  ikke et Blink med Øiet
%          p. 28  „störe Gud“ — a WRONG SORT: store is meant, and the same
%                 Kierkegaard sentence is set „store Gud“ with a plain o on
%                 p. 24. Logged below as a printer's error as well.
%        pp. 1–15 contain no ø at all.
%
%  SECOND QUOTATION PAIR. On pp. 6–7 the fount runs out of the proper sorts
%        and the compositor sets five quotations with substitute ,, and ''
%        instead of „ and “. These are set here as „ … “ and each site is
%        flagged in a % comment — the one place in this file where a
%        typographic accident has been normalised rather than reproduced,
%        because the substitute sorts carry no sense and the „…“ pair is what
%        the compositor was reaching for. Reversible from the comments.
%
%  PRINTER'S DEFECT LOG (all transcribed as printed, flagged in % comments at
%  their sites):
%        p. 1   words dropped: „er jeg saa rigtignok ogsaa Magt“
%        p. 10  „charakterisk“ for charakteristisk
%        p. 13  „i“ set twice across the line break
%        p. 14  battered em-dash after „Troen“
%        p. 22  „christendum“ — the word-space dropped
%        p. 26  broken em-dash sort
%        p. 27  two damaged em-dashes (one broken at the waist, one printing a
%               vertical shoulder)
%        p. 28  „störe Gud“ for store Gud (see the ö/ø census above)
%        p. 29  „först“ — the t has lost its left crossbar and prints like f/ſ
%        p. 29  „Forudsætninger.,“ — a full point set before the comma
%        p. 29  „komme“ for komne
%        Battered and part-filled t and r sorts are common throughout and are
%        not individually logged. Ink specks that could be read as accents
%        (p. 6 „et“, p. 11 „Side“, p. 21 before „den bagvendte“) are read as
%        specks and not transcribed.
%
%  VARIANTS WITHIN THE ARTICLE, recorded and not normalised:
%        p. 18  „min Persons Forskjellighed“ where p. 15 has Personligheds
%        p. 21  „Inderligjörelse“   vs   p. 22  „Inderliggjörelse“
%        „Til Selvprövelse“ is letterspaced on p. 18 but NOT on pp. 21 and 27.
%
%  ELLIPSES, as counted on the images: five points on pp. 1, 16 (×3) and 17;
%        four on pp. 19, 21, 22 (×3), 23 (×3), 24, 25 (×2), 26; three on
%        pp. 11, 26 and 27.
%
%  QUOTE BALANCE: +4, and all four are deliberate. This book's nested
%        quotations REUSE the „…“ pair rather than switching sorts, so an
%        opening mark is regularly left unclosed where the inner quotation
%        ends. The four sites: p. 12 „Hvor er i Guds Huus… (closes on an
%        ellipsis, no closing mark); p. 22 „Skulde dette höre…; p. 23 „Lad os
%        antage…; p. 25 the letterspaced display. Nobody should "fix" this to
%        zero. (p. 7 has an orphan CLOSING mark after „Öine,“ which cancels an
%        earlier opening, so the raw count understates the reuse.)
%
%  EMPHASIS: 12 calls, all letterspacing. pp. 1 (Scharling, Engelstoft),
%        2 (the thirteen-line Kierkegaard quotation), 11 (the title-list),
%        12 (Christelige Talers; „Forvar din Fod…“), 14 („Til Selvprövelse“,
%        „Indövelse i Christendom“), 17 (the display from „Var“ through
%        „— er dette Sandhed?“, resuming plain at „vare der vistnok“),
%        18 („Til Selvprövelse“), 25 (the display „Hvad der skal gjöres…“,
%        ending after „Anden.“; „Exempel:“ and the bracketed source are plain),
%        25 (in the footnote, „Aarstal og Datum“).
%
%  TERMINAL RULE on p. 29 measures 0.26 of the measure; set at 0.25.
% ============================================================
\documentclass[12pt,a4paper]{article}
\usepackage[utf8]{inputenc}
\usepackage[T1]{fontenc}
\usepackage[danish]{babel}
\usepackage[protrusion=true,expansion=true]{microtype}
\usepackage[margin=1.3in]{geometry}
\usepackage{setspace}
\usepackage{libertinus}
\usepackage{libertinust1math}
\usepackage{fancyhdr}
\usepackage[colorlinks=true,linkcolor=black,urlcolor=black,
  pdftitle={Om S. Kierkegaards „mentale Tilstand“},
  pdfauthor={R. Nielsen}]{hyperref}

\setlength{\headheight}{15pt}
\pagestyle{fancy}
\fancyhf{}
\fancyhead[L]{\textit{Om S.\ Kierkegaards „mentale Tilstand“}}
\fancyhead[R]{\thepage}
\setlength{\parindent}{1.5em}
\setstretch{1.3}

% the book marks both its footnotes *) and restarts per printed page
\renewcommand{\thefootnote}{*)}

\begin{document}

% --- p. 1 ---
% Head of the article. Folio AND running head are both suppressed on this page;
% the sheet signature "1" at the foot right is not transcribed.
% Title and byline are centred in the same antiqua, larger and heavier than the
% body; a short centred rule stands above and below the epigraph.
\begin{center}
  {\large Om S.\ Kierkegaards „mentale Tilstand“.}

  \medskip
  Af Prof.\ R.\ Nielsen.
\end{center}

\bigskip
\begin{center}\rule{0.25\linewidth}{0.4pt}\end{center}
\bigskip

\begin{quote}\small
% EPIGRAPH, quoted from Nyt Theologisk Tidsskrift 8:1 (1857), 203--204.
% It opens with FIVE spaced full stops and then a large two-line raised initial
% T (of „Til“); the initial is not reproduced here.
% LETTERSPACING: „Scharling“ and „Engelstoft“ are letterspaced — confirmed on
% the image at 600 ppi. Nothing else in the epigraph is.
% ö/ø: hörte, sögte, Berömmelse are all o-DIAERESIS. No ø on this page.
. . . . . Til noget af det Besynderligste, som vi have oplevet i vore
Dage, hörte ikke saameget det, at et Menneske, om hvis mentale
Tilstand i det sidste Afsnit af hans Liv, man ikke kan have nogen
sikker Dom, sögte en herostratisk Berömmelse ved at udslynge For\-%
nærmelser og Beskyldninger mod en Mands Charakteer . . . ., men
% four points before the comma here; five after "var." below. Counted at 1200 ppi.
endnu forunderligere var. . . . .\quad Nyt Theologisk Tidsskrift, udgivet
af Dr.\ C.\ E.\ \emph{Scharling} og Dr.\ C.\ T.\ \emph{Engelstoft}, 8de Bind.
1ste Hefte 1857, S.\ 203 og 204.
\end{quote}

\bigskip
\begin{center}\rule{0.25\linewidth}{0.4pt}\end{center}
\bigskip

For at kunne have „nogen sikker Dom“ om S.\ Kier\-%
kegaards „mentale Tilstand“, være sig i det förste eller
det sidste Afsnit af hans Liv, er det en ufravigelig Be\-%
tingelse, at man dog med nogen Eftertanke besinder sig
paa, hvad det da i sin Heelhed er, „dette Menneske“
egentlig vilde, hvilke Kræfter han udfoldede, hvilke Mid\-%
ler han anvendte, hvilke Vaaben han brugte.

„I Qvaler, som sjældent et Menneske oplevede dem,
i Aandsanstrængelser, der i Löbet af 8 Dage vel vilde
beröve en Anden Forstanden, er jeg saa rigtignok ogsaa
% PRINTER'S ERROR: words have dropped out here. „er jeg saa rigtignok ogsaa
% Magt“ is unintelligible; Kierkegaard's own text runs „...er jeg saa rigtignok
% ogsaa i Besiddelse af en Magt“. Verified at 900 ppi that nothing stands at
% either line end. Transcribed as printed.
Magt, unegtelig en for et stakkels Menneske forförerisk
Bevidsthed, dersom Qvalen og Anstrengelsen ikke i den
Grad var det Overveiende, at ofte nok mit Önske er
% --- p. 2 ---
% The pencilled shelf-mark "Scan 6664.585.5", the marginal pencil note beside
% the letterspaced quotation, and the Harvard date-stamp are not transcribed.
Döden, min Længsel Graven, min Begjæring, at mit Ön\-%
ske og min Længsel snart maatte opfyldes.“

Den, der saaledes selv anmelder sin Optræden, maa,
dersom han ellers er ved Samling, tilvisse have en Byrde
paa sin Samvittighed; maa, dersom han virkelig veed,
hvad han siger, for Alvor have Noget at sige. Hvad
vil han da? „Redelighed!“

Og hvad er det saa han har at sige?

% LETTERSPACED (Sperrsatz) throughout, from „Dette to stendom.“ — thirteen
% lines. The compositor closes up the last words of the fullest lines
% („har Du“, „mindre, Du“), but the run is continuous; checked line by line
% on the image. This is the long letterspaced quotation of the article.
\emph{„Dette skal siges, saa være det da sagt:
Hvo Du end er, hvilket Dit Liv forresten, min
Ven, --- ved (hvis Du ellers deeltager i den) at
lade være at deeltage i den offentlige Guds\-%
dyrkelse som den nu er (med Paastand paa at
være det nye Testamentes Christendom) har Du
bestandigt een og en stor Skyld mindre, Du
deeltager ikke i at holde Gud for Nar ved
at kalde det nye Testamentes Christendom,
hvad der ikke er det nye Testamentes Chri\-%
stendom.“}

Det er stærke Ord, en Gru opvækkende Dom. Er
Den, der tör tale saaledes, og dömme saaledes, nu ogsaa
virkelig en Berettiget; trues da ikke Samfundet med en
Aandskrise, der i sine Fölger kunde blive betænkelig for
Staten, farlig for Folkets aandelige Ro?

„Hermed --- ja, o Gud skee nu hvad Din Villie er,
uendelige Kjærlighed! --- hermed har jeg talet.“

Dette er jo vistnok religiöst; men er det dog maa\-%
skee ikke igjen altfor religiöst?

„Maaskee skal Een höre saaledes, at han gjör, hvad
jeg siger, en Anden saaledes, at han forstaaer det som
Gud velbehageligt, mener at vise Gud en Dyrkelse ved
at deeltage i at löfte Skrig imod mig: ingen af Delene
vedkommer mig; mig vedkommer kun, at det skal siges.“

Her er i alt Dette en Resignation, en Besluttethed,
en Pathos, der unegtelig tyder hen paa en ganske sær\-%
egen „mental Tilstand“; men hvilken? Man lytter, og
% --- p. 3 ---
% The sheet signature "1*" at the foot right is not transcribed.
% ö/ø: spörger, Öjeblik, Tröst, velgjörende, indrömme, gjöre, paanöde, förste
% (twice), mörke, Indrömmelse — every one o-DIAERESIS at 900--1200 ppi. No ø.
spörger med Ængstelse hvilken? Lykkes det nu i et saa
kritisk Öjeblik Tidens meest Klarsynede at gjennemskue
den Optrædendes „mentale Tilstand“, og overbevise sig
om, at det ingenlunde er en Berettiget, man har for sig,
men en sindssvag Herostrat, et stakkels Genie, der i
Overspændthed er gaaet saavidt, at han er gaaet fra
Forstanden; da fatter Samtiden strax, at Manden er gal, og
% a speck of ink above the comma after „Tröst“ makes tesseract read a
% semicolon; at 900 ppi it is a comma.
det er Samtiden en stor Tröst, „at Manden er gal“. Thi, det
forstaaer sig, naar Manden er gal, saa er Spændingen
endt, Krisen overstaaet, og „Redeligheden“ falder af
sig selv.

De Klarsynede, som speculere i Fordelen af at ud\-%
brede og vedligholde den for Sjælenes Ro saa velgjörende
Mistanke mod Ks.\ „mentale Tilstand“, ere da naturlig\-%
viis ogsaa paa det Rene med sig selv om, hvorledes
denne Mistanke kan fra mere end eet Synspunkt finde
sin Bestyrkelse.

Man kan fra den ene Side indrömme Sandheden af
Alt, hvad K.\ siger om Forskjellen imellem det nye Te\-%
stamentes Christendom og Nutidens Christendom; man
kan beundre den Skarpsindighed, Klarhed, Conseqvens
og Sikkerhed, hvormed han forstaaer at gjöre Uligheden
indlysende, og endda finde hans „mentale Tilstand“ be\-%
klagelig. Hans Aandsforvirring bestaaer da ingenlunde
% an ink speck sits under the second e of „seer“; the word is „seer“.
deri, at han ikke seer, hvad Christendom oprindelig er,
men deri, at han for ramme Alvor vil paanöde os „det
nye Testamentes Christendom“, vil skrue Tidsalderen til\-%
bage til hiint historisk overvundne Standpunkt; „altsaa
endnu hiin förste ubetingede Tro paa Gudmenneskets
Mirakler, endnu hiint förste asketiske Verdenshad, hiin
mörke Livsanskuelse: Manden er jo gal!“

Fra den anden Side begynder man derimod med
den Indrömmelse, at det nye Testamentes Christendom
er den eneste for alle Tider gjældende Christendom.
Ks.\ Feil er da ikke den, at han endnu i vore Dage
fordrer Evangeliet anerkjendt; men Vildfarelsen er, at
% --- p. 4 ---
% ö/ø: nödvendigste, tör, anföres (twice), nödvendigt, förste (twice), Höide,
% höist, höi, frastöde — all o-DIAERESIS, checked at 900 ppi. No ø.
han ikke ret kjender Evangeliet, efterdi han ikke seer,
at det netop er det nye Testamentes Christendom, som
for Tiden finder sit væsentligste Udtryk, sine bedste Be\-%
tingelser og nödvendigste Garantier i det Bestaaende.
„Ulykkelige Sindssvaghed, at han dog ikke kunde see,
at det kun er et Udvortes, her er forandret, ikke det
Indvortes; ikke see, at Menighedslivet endnu er beteg\-%
net med den samme Broderkjærlighed, som i det nye
Testamente; ikke fik Naade til at see, at Ordets Tjenere,
dersom det kom derpaa an, strax ere villige til at for\-%
lade Alt for Christi Skyld, og altsaa endnu, paa nogle
Smaaforandringer nær, ligne Disciplene i det nye Testa\-%
mente; den stakkels Forblindede!“

Den Omstændighed, at man saaledes ad to forskjel\-%
lige Veie kan komme til samme Resultat angaaende Ks.
„mentale Tilstand“, tör man imidlertid ikke forhaste sig
med at udtyde, som havde vi her to Beviser, der smel\-%
tede sammen til gjensidig Forstærkning. De angivne
Synspunkter ere jo hinanden saa modstridende, at hvad
der fra den ene Side anföres som Beviis paa Sanddruhed,
maa fra den anden Side anföres som Tegn paa Afsindig\-%
hed, en indre Modsigelse, der just ikke egner sig til at
sikkre Dommen.

Nei, for at have en sikker Dom om Ks.\ „mentale
Tilstand“ i det sidste Afsnit af hans Liv, er det nödven\-%
digt, at man forud danner sig en sikker Dom om hans
„Forfattervirksomhed“ i det förste Afsnit. Denne Dom
har i selvbehagelig Overfladiskhed hidtil holdt sig saa
svævende, og heri ligger Grunden til, at Misforstaaelsen
paa det Sidste maatte voxe til en saadan Höide.

Vistnok vakte K.\ ved sin förste Optræden Opmærk\-%
somhed, men det i hans Stiil og Skrivemaade saa höist
Eiendommelige bevirkede strax, at han i ligesaa höi
Grad maatte frastöde Nogle, som tiltrække Andre. Til
et andet, for dybere gaaende Kritik mere gunstigt Tids\-%
punkt kunde dette maaskee have havt den gavnlige
% --- p. 5 ---
% ö/ø: Fölge, Först, fölge, Skjön, Religiöse, förste, föler, Rörende, afdöde,
% Sönnen (twice), böie, hört, fölte, skjönt — all o-DIAERESIS. No ø.
Fölge, at hans Anskuelse strax fra de modsatte Stand\-%
punkter var bleven saa stærkt belyst, at den almindelige
Bevidsthed fra Först af havde kunnet fölge med; nu
derimod blev Opfatningen i det Hele overladt det indi\-%
viduelle Skjön, Kritiken gik om i Negligé: de forskjellige
Meninger gjaldt efter Smag og Behag.

Saaledes, for her at holde os til det Religiöse, mente
Nogle, at K.\ egentlig manglede Alvor, medens Andre
fandt hos ham den strengeste Alvor. Blandt disse Sidste
stod endog Biskop Mynster selv. Allerede de förste
„opbyggelige Taler“ gjorde et saa alvorligt Indtryk paa
den gamle Biskop, at denne, idet han med Vægt udtaler
sin Anerkjendelse, endog föler sig bevæget til at give et
psychologisk Bidrag til nærmere Forstaaelse af den unge
Mands besynderlige Alvor.

„Det har for mig,“ siger Biskop Mynster (Kirkelig Pole\-%
mik. Intelligensblade, udg.\ af J.\ L.\ Heiberg. Jan.\ 1844.)
noget Rörende, at Mag.\ Kierkegaard stedse helliger sine
opbyggelige Taler til sin afdöde Faders Minde. Thi og\-%
saa jeg har kjendt denne agtværdige Mand; han var en
slet og ret Borger, gik fordringsfri sin stille Gang gjen\-%
nem Livet, havde aldrig gjennemgaaet nogetsomhelst
„philosophisk Bad“. Hvorledes skeer det da, at Sönnen
i sin rige Dannelse, saa ofte han vil nedskrive opbygge\-%
lige Taler, stedse vender Tanken hen til hiin Mand, der
længesiden er indgangen til Hvile? Hvo der har læst
den deilige Tale --- eller lad os kun kalde den Prædi\-%
ken --- „Herren gav, Herren tog, Herrens Navn være
% the same Job text is quoted again on p. 9, there with a semicolon:
% „...Herrens Navn være lovet;“. Both readings stand as printed.
lovet,“ vil forstaae det. Sönnen har, ligesom jeg, seet
den gamle Fader ved bittre Tab, han har seet ham folde
sine Hænder, böie sit ærværdige Hoved, han har hört
hans Læbe udsige hiint Ord, men tillige seet hans hele
Væsen udsige det paa en saadan Maade, at han fölte,
hvad han saa skjönt udvikler om Job, at „ogsaa den er
en Menneskenes Lærer, der ingen Lære havde at over\-%
give Andre, men efterlod Slægten sig selv som et For\-%
% --- p. 6 ---
% ö/ø: fölt, Sönnen, Sörgehuset, fölende, religiöse (twice), religiös,
% frastödte (twice), höieste, Öine, berömte, tilföier — all o-DIAERESIS. No ø.
billede, sit Liv som Veiledning for ethvert Menneske“,
der saae det; han har fölt, at den Gamle havde i sit
fromme Ord overvundet Verden. Og hvad Sönnen der
i Sörgehuset saae hos sin gamle Fader, det nedskrev han
i en Prædiken, som vil tiltale, vederqvæge ethvert fölende
Hjerte.“ . . . .

Var den almindelige Dom om Ks.\ „religiöse Alvor“
allerede noget usikker, saa blev den videnskabelige Dom
om hans religiöse Begrebsudvikling endog aldeles for\-%
% a small slanted ink mark stands over the e of „et“ here, and both OCR
% witnesses read „ét“. Set „et“: the same phrase is printed „Troen er et
% Paradox“, unaccented, on p. 10, and stray specks of this size are common
% on these pages (cf. „Side“ p. 11, „seer“ p. 3).
virret; han kaldte jo Troen et Paradox, og betjente sig
% QUOTE SORTS: from here to the end of p. 7 the compositor twice sets the
% quotation marks from SUBSTITUTE sorts — two separate low commas for the
% opening and two raised turned commas for the closing — visibly wider and
% lighter than the „ and “ sorts used everywhere else. Sites: „den indirecte
% Meddelelsesform“ and „Paradoxet“ here, „höieste Problemer“ and „Frygt og
% Bæven“ on p. 7, and the closing after „Öine,“ on p. 7. Normalised to „…“
% here; the substitution is recorded, not set.
af „den indirecte Meddelelsesform“: han frastödte de
Speculative, han frastödte Theologerne.

At Forkastelsen af den indirecte Form dog endnu
ligesaalidt var eenstemmig, som Forargelsen paa „Para\-%
doxet“ almindelig, kan udtrykkelig sees blandt Andet af
den Maade, hvorpaa Biskop Mynster i ovennævnte Ar\-%
tikel tillige omtaler „Frygt og Bæven“. „Jeg har ogsaa“,
siger Biskoppen, „læst det mærkelige Skrift: „Frygt og
Bæven“, og hvad jeg endog kan have savnet deri, da
visselig ikke en dyb religiös Grund, ikke et Sind, der
er mægtigt til at træde Livets höieste Problemer under
Öine; det har levende mindet mig om det berömte Sted
hos Jacobi: „Ja, ich bin der Atheist und Gottlose, der
dem Willen, der Nichts will, zuwider, lügen will, wie Des\-%
demona sterbend log, lügen und betrügen will, wie der
für Orest sich darstellende Pylades u.\ s.\ w.\ (Jacobi an
Fichte S.\ 32), hvoraf det dog ingenlunde er en Efterlig\-%
ning eller Efterklang.“ Under denne Vending lægger
altsaa Biskoppen selv Vægt paa den indirecte Meddelelse,
og afviser dem, der vilde beskylde Forfatteren for „uti\-%
dig Jagen efter det Pikante“, idet han, efter at have
udtalt sig angaaende Bogens Titel, udtrykkelig tilföier:
„Saaledes har denne Kamp under Frygt og Bæven ikke
det Mindste tilfælles med den velbehagelige Leflen, hvor\-%
med de moralske Genier hige efter at komme i saadanne
Forhold, hvor de kunne vise, at de ikke forgjæves have
% --- p. 7 ---
% The sheet signature "2" at the foot right is not transcribed.
% ö/ø: göttlich, förste, mödte, religiös (three times), höieste (twice), Öine
% (three times), stöttede, indrömmer, Selvfölge, indrömme, nödvendige,
% Indrömmelsen, hörer, Spög — all o-DIAERESIS. No ø.
læst Schelling, at de ogsaa ere istand til „einzig göttlich
zu handlen.“

Havde Biskop Mynster nu virkelig kunnet bevare
dette sit förste Indtryk af, at han her i „Frygt og Bæ\-%
ven“ mödte „en dyb religiös Grund“, et Sind, der var
mægtigt til at træde Livets höieste Problemer under
% UNBALANCED QUOTATION, as printed: this closing mark (in the substitute
% sort) has no opening mark of its own — Nielsen runs Mynster's words on
% from the quotation of p. 6. Logged, not repaired; nested quotations in this
% article reuse the same pair, so the balance count is unreliable here.
Öine,“ havde han i denne Henseende havt Tillid til sin
egen Dom, var selv traadt disse „höieste Problemer“
under Öine, og saaledes ved sit Exempel bleven en Vei\-%
ledning for de Mange, der saa gjerne stöttede sig til
hans Autoritet: hvor vilde Sagens Gang da ikke fra Be\-%
gyndelsen af have taget en ganske anden Vending!
Naar man indrömmer en Forfatter hans Forudsætninger,
da er det jo en Selvfölge, at man ogsaa maa indrömme
ham Forudsætningernes nödvendige Conseqvenser, at sige,
dersom Indrömmelsen beroer paa Indsigt, og ikke paa en
flygtig Stemning. Men i „Frygt og Bæven“ have vi jo
de principielle Forudsætninger. Det hörer just med til
det Eiendommelige i Kierkegaards Væsen, at han fra
Begyndelsen af er saa enig med sig selv, saa indvortes
færdig. Han var ikke en Ung, der med Aarene blev
gammel, ikke en Lystig, der siden blev saa alvorlig,
ikke en Æsthetiker, der senere blev religiös; nei,
han var oprindelig Alt, hvad han var, i eiendom\-%
melig Fordobling, i sin Ungdom gammel, i sin Spög al\-%
% spacing.py flags a run of spaced punctuation in the next line
% („glad , i sin Strenghed mild , i sin“); it is justification of a stretched
% line, not copy, and is closed up.
vorlig, i sin Smerte glad, i sin Strenghed mild, i sin
Bitterhed veemodig: Kierkegaard er i den Grad en aprio\-%
risk Natur, at han næsten mangler det Perfectible.

I „Frygt og Bæven“ er Principet udtalt, fuldstændig ud\-%
talt; vi see i dette Skrift som i et Speil Forfatterens hele in\-%
tellectuelle Physiognomie. Er der nu i dette Physiognomie
noget væsentlig Fordreiet; da er det i Sandhed paafaldende,
at en Mand, som Biskop Mynster, har gjennem disse for\-%
dreiede Træk kunnet mene at see --- „en dyb religiös
Grund“. Eller, dersom her i dette Anlæg er noget Falsk,
og dette Falske, hvad det jo maa være, forsaavidt det
% --- p. 8 ---
% ö/ø: Öie, höieste, Öine, Rastlös, overströmmende, Slör, gjennemföre,
% höieste, löst, förste, hörte, frastöde, föiede, nöd, Strötanker,
% iöinefaldende — all o-DIAERESIS. No ø.
ligger i Forudsætningerne og fremgaaer af Principerne,
endog er et Grundfalsk; er det da ikke uforklarligt, at
en Kjender som Biskop Mynster, istedetfor strax at faae
Öie paa dette Grundfalske, har kunnet lade sig blænde,
endog i den Grad, at han bag den pseudonyme Charac\-%
teermaske har troet at opdage en med sin egen beslæg\-%
tet Aand, „et Sind, der var mægtigt til at træde Livets
höieste Problemer under Öine“?

Rastlös og med overströmmende Productivitet arbei\-%
dede K.\ i en Række Aar paa at klare, hvad han kaldte
sine „Existensproblemer.“ Hans indirecte Meddelelse,
der bag sit Slör bevarede Inderligheden, mildnede det
Polemiske, og satte forligende Skilsmisse imellem Virkelig\-%
heden og Idealiteten, skaffede Ro, saa at den taktiske
Fremrykken kunde skee i al Stilhed.

Det var Kierkegaards Plan at gjennemföre en Re\-%
vision af Existenskategorierne, og udskille dem af den
Forplumring, hvori en misledet Videnskabelighed havde
neddraget „Livets höieste Problemer“. Han har löst
denne Opgave med en Sikkerhed, Grundighed og Alsi\-%
dighed, der i kommende Tider vil hjemle ham Plads
blandt Tænkerne af förste Rang. Men trods hans Ag\-%
telse for sand Videnskabelighed, hvor Videnskab gjælder,
og uagtet han i Gjerningen er Systematiker trods Nogen;
saa hörte det dog væsentlig med til den Skjærpelse af Inder\-%
lighedsproblemerne, han havde foresat sig, idelig, endog
systematisk, at bryde den systematiske Form, overalt at fra\-%
stöde det dumt Videnskabelige. Han anvendte denne Taktik
med stor Behændighed, föiede Bygning til Bygning, og nöd
omsider den Tilfredsstillelse, at Krandsen allerede vaiede
paa Anlægets sidste Bygning, da paragraphisk Alvor kom
til, og naturligviis ansaae det Hele for en Masse „Strö\-%
tanker, Aphorismer, Indfald og Glimt.“

At Kierkegaard vilde „Speculationen“ tillivs, og at
det især gik ud over „Systemet“ var en saa iöinefaldende
Kjendsgjerning, at der endog var dem, som mente, at
% --- p. 9 ---
% ö/ø: dömmer, Jödedom, ifölge, religiöst-ethiske, nödvendigt, gjör, Fölelse,
% Fölgende, hörer, berört, fölende — all o-DIAERESIS. No ø.
han udgav Troen for et Paradox alene for at drille „Sy\-%
stemet“. Herom kan man kun sige, at Enhver dömmer,
som han har Forstand til, og derhos minde om „Stadier\-%
nes“ Motto: „Solche Werke sind Spiegel, wenn ein Affe
hinein guckt, kann kein Apostel heraus sehen.“

Hvor magtpaaliggende det maatte være K.\ reent
praktisk og for Existensens Skyld at faae „Speculatio\-%
nen“ adskilt fra „Troen“, vil man forstaae, naar man
seer hen til den dogmatisk-speculative Retnings inderlige
Sammenhæng med, hvad den nyere Philosophie fra Spi\-%
noza af har kaldet „Immanensen“. I Immanensen er
ikke alene den hedenske Ethik i dens ædleste, sokratiske
Form, begrundet, men selv indenfor Jödedom og Chri\-%
stendom gives der, ifölge Ks.\ Opfatning, en folkelig hu\-%
man Tilegnelse af den aabenbarede Religion, der i en
vis Maade ligeledes svarer til Immanensen. Det er saa
langt fra, at K.\ vil have dette religiöst-ethiske Mellem\-%
stadium overseet, forkastet eller ringeagtet, at han tvert\-%
imod har stræbt at udvikle dets Betydning i en Række
af „opbyggelige Taler.“

Til „en sikker Dom“ om Ks.\ Forfattervirksomhed
er det nödvendigt, at her igjen lægges Eftertryk paa
denne, under de seneste Stridigheder kun altfor meget
overseete, Side af hans Productivitet; thi disse Taler have
just det Milde, Besindige, Stilfærdige, det Vækkende og
dog Beroligende, som man ellers mener at savne hos
K., og hvori man saa pleier at sætte det ægte Evange\-%
liske. Hvilken Forskjel det i saa Henseende gjör, om
man i sin Dom ledes af umiddelbar Fölelse, eller af en
tydelig Tanke, en bestemmende Kategorie, kan sees
af Fölgende. Blandt de opbyggelige Taler hörer just
ogsaa den over Jobs Ord: „Herren gav, Herren tog,
Herrens Navn være lovet;“ denne Tale fandt Biskop
Mynster, som allerede berört, saa opbyggelig, at den
maatte kunne „tiltale, vederqvæge ethvert fölende Hjerte“,
og vilde saa af den Grund kalde den en Prædiken. Her
% --- p. 10 ---
% ö/ø: berömte, fölende, Höide, gjör — all o-DIAERESIS. No ø.
% PRINTER'S ERROR in the first word-group: „charakterisk“ for
% „charakteristisk“. Confirmed at 1200 ppi; both OCR witnesses agree.
% Transcribed as printed.
er et charakterisk Træk. Den berömte Prædikant ap\-%
pellerer til „ethvert fölende Hjerte“, og vil have den
opbyggelige Tale om Job kaldet „Prædiken“, hvorved
han da tilkjendegiver, at den i Talen udviklede Livsan\-%
skuelse holder Höide med den i det kirkelige Foredrag;
den dialektiske Forfatter, hvis Bestræbelse det er, at faae
de forskjelligartede Existenskategorier udklarede, vil
derimod ikke have sin Tale kaldet Prædiken, fordi han
vil, at „Prædiken“ skal svare til det Christelige, og veed
med sig selv, at denne hans Tale ikke udtrykker det i
egentlig Forstand Christelige. Hans Replik igjennem
Pseudonymen i „Uvidenskabelig Efterskrift“ S.\ 204 og
205 viser dette. „Mag.\ Kierkegaard har formodentlig nok
vidst, hvad han gjorde, da han kaldte de opbyggelige
Taler: opbyggelige Taler, samt hvorfor han afholdt sig
fra Brugen af christelig-dogmatiske Bestemmelser, fra at
nævne Christi Navn o.\ s.\ v., hvad man ellers i vor Tid
gjör frisk væk, medens Kategorierne, Tankerne, det Dia\-%
lektiske i Fremstillingen kun er Immanensens. Som de
pseudonyme Skrifter, foruden hvad de directe ere, indi\-%
recte ere en Polemik mod Speculationen, saaledes ere
disse Taler det ogsaa, ikke derved, at de ikke ere spe\-%
culative, men derved, at de ikke ere Prædikener.“

Man vil let forstaae, at her ikke er Tale om „en
christelig Strenghed ligeoverfor en given christelig Mild\-%
hed“; her handles om at finde „Redelighed“ i Principerne,
hævde Evangeliet dets Eiendommelighed, og afværge,
at det kategorisk forvexles med et dermed beslægtet,
men dog underordnet Resignationens og Fromhedens
Standpunkt. De, som med nogen Opmærksomhed og Ef\-%
tertanke have fulgt den kierkegaardske Bevægelse, maae
jo vide, at denne Forfatter ikke lod det beroe ved at
lade nogle Pseudonymer opbyde en glimrende Dialektik
for at udvise „Stadier paa Livets Vei“, og med formel
Skarpsindighed hævde den Paastand, at Troen er et Pa\-%
radox; men, at det i den Grad var ham om Sagens rette
% --- p. 11 ---
% ö/ø: religiös-praktiske, Gjennemförelse, gjöre, religiöst, Angerlöse —
% all o-DIAERESIS. No ø.
Forstaaelse og religiös-praktiske Gjennemförelse at gjöre,
at han i en rig Mangfoldighed af populære Foredrag:
% LETTERSPACED: the run of Kierkegaard's titles, from „Opbyggelige“ through
% „Fredagen“. The following „o. s. v.“ is set in ordinary spacing.
\emph{Opbyggelige Taler i forskjellig Aand, Kjærlighe\-%
dens Gjerninger, Christelige Taler, Taler ved
Altergangen om Fredagen} o.\ s.\ v., har paa det
Omhyggeligste viist, og i de forskjelligste Fortoninger
belyst, hvad Kategorierne gjælde, hvorledes der, om
muligt, skulde virkes paa Folket, hvorledes der skulde
tales i Menigheden.

„Menigheden! Men er det da virkelig ikke sandt,
hvad dog fra saamange Sider forsikkres, at S.\ Kierke\-%
gaard kun vilde være „hin Enkelte“, og forsaae sig paa
„den Enkelte“, i den Grad, at han aldeles oversaae
Menigheden?“ Nei, det er virkelig usandt, kategorisk
usandt. Kun ved forstærket Eftertryk paa „den Enkelte“
blev det ham muligt at sikkre Menighedsbegrebet mod
Forvexling med Samfund af lavere Orden; „forsaavidt“,
siger han, „der, religiöst, er „Menighed“, da er dette et
% a thin slanted ink mark stands over the final e of „Side“ here. Read as a
% speck, not an acute: „Sidé“ is not a word, and specks of this size recur on
% these pages (cf. „et“ p. 6, „seer“ p. 3). Set „Side“.
Begreb, som ligger paa den anden Side af „den Enkelte“,
og som for Alt ikke maa forvexles med, hvad der poli\-%
tisk kan have Gyldighed, Publikum, Mængde, det Nu\-%
meriske o.\ s.\ v.“

% THREE spaced full stops here, counted at 900 ppi — the only three-point
% ellipsis in pp. 1--15; the others run to four or five.
„Mig, den opbyggelige Forfatter, var og er det der\-%
for en Glæde, at ogsaa . . . De blive flere, som blive
opmærksomme paa det om „den Enkelte“; det var og
er mig en Glæde, thi vistnok har jeg Tro paa min Tankes
Rigtighed trods hele Verden, men hvad jeg dog næst
den sidst opgav, er min Tro paa de enkelte Mennesker.
Og dette er min Tro, at saa meget Forvirret og Ondt
og Afskyeligt der kan være i Menneskene, saasnart de
blive det Ansvar- og Angerlöse „Publikum“, „Mængde“
o.\ D., ligesaa meget Sandt og Godt og Elskværdigt er
der i dem, naar man kan faae dem Enkelte. O, og i
hvilken Grad skulde Menneskene ikke blive --- Mennesker
og Elskelige, naar de vilde blive Enkelte for Gud!“
(„Om min Forfattervirksomhed“, S.\ 12).
% --- p. 12 ---
% ö/ø: Religiöse, spögende, Röst, Söndagspræst, höres, förste, Höiröstethed,
% gjöre (twice), höie, Skjönhed, stöde — all o-DIAERESIS. No ø.
Den udelukkende Interesse, hvormed K.\ arbeidede
for det Religiöse, gav sig da ogsaa tilkjende derved, at
han ikke blot skrev Prædikener, men tillige nu og da
paa given Anledning virkelig prædikede, helst til Alter\-%
gang om Fredagen. Hvorfor han just foretrak Fredagen,
forklarede han engang i en Samtale paa sin spögende
Maade, idet han beklagede, at han ikke havde Röst
til at være Söndagspræst, og höres af Saamange, men
dog mente, at han med sin svagere Stemme vel kunde
være skikket til Fredagspræst.

Til nærmere Bestemmelse af det Synspunkt, hvor\-%
under K.\ vilde have det i egentlig Forstand christelige
% LETTERSPACED: the book-title „Christelige Talers“ and, below, the whole of
% the sermon-title „Forvar din Fod, naar Du gaaer til Herrens Huus.“
Foredrag opfattet, ville vi her henvise til \emph{Christelige
Talers} tredie Afdeling, og fremhæve nogle Punkter af
den förste Tale: „\emph{Forvar din Fod, naar Du gaaer
til Herrens Huus.}“ Talen begynder saaledes:

„Hvor er i Guds Huus Alt saa stille, saa trygt.
Den, der træder derind, ham er det, som var han med
et eneste Skridt kommen til et fjernt Sted, uendelig langt
borte fra al Larm og Skrig og Höiröstethed, fra Tilvæ\-%
relsens Forfærdelser, fra Livets Storme, fra rædsomme
Begivenheders Optrin eller deres forsmægtende Forvent\-%
ning. Og hvor Du vender dit Blik hen derinde, vil Alt
gjöre Dig tryg og rolig. Den ærværdige Bygnings höie
Mure, de staae saa faste, de værne saa tilforladeligt om
det sikkre Tilflugtssted, under hvis mægtige Hvælving Du
er fri for ethvert Tryk. Og Omgivelsens Skjönhed, dens
Pragt vil gjöre Dig Alt venligt, saa indbydende, den vil
% the quotation closes on an ellipsis with NO closing quotation mark; as
% printed, not repaired.
ligesom indynde det hellige Sted hos Dig. . . . .

Læseren forstaaer, at denne tiltrækkende Beskrivelse
netop skal bruges til dermed at stöde fra. „Hvor bero\-%
ligende, hvor inddyssende, ak, og hvormegen Fare i
denne Tryghed!“

Ja, men Taleren er alligevel ingen Fanatiker, der
kun veed at ivre mod et Udvortes, ingen Billedstormer,
der reformerer ved at slaae istykker; heller ikke er
% --- p. 13 ---
% ö/ø: söge, höie, Oplöftelse, religiös, höre, Prövelser, försögt, Öiemed,
% Höide, nödvendigviis, söge, lös, religiöse, Religiöse — all o-DIAERESIS.
% No ø.
han en Ukjærlig, der ingen Sympathi kan have med
dem, som fra „Larm og Skrig“, fra „Tilværelsens For\-%
færdelser“ og fra „Livets Storme“ söge et „Tilflugtssted“
bag de „höie Mure“, en Sjælens Oplöftelse under den
„mægtige Hvælving“. Det er en religiös Tænker, der i
Talen vil advare mod den farligste, meest bedaarende
af alle Illusioner, Andagtens nemlig, naar Andagten er
en æsthetisk Stemning, en blid Illusion. „Hvor beroli\-%
gende, hvor inddyssende, ak, og hvormegen Fare i
denne Tryghed! Og derfor er det dog sandeligen saa,
at det egentligen kun er Gud i Himlene, der i Livets
Virkelighed ret tilgavns kan prædike for Menneskene,
thi han har Omstændighederne, har Skjæbnerne, har
% four points after „Magt.“, five after „Beroligelse!“ — counted at 900 ppi.
Bestræbelserne i sin Magt. . . . . Men i Guds Huus, i
% „i Guds Huus, i / i det pragtfulde Guds Huus“ — the „i“ is printed twice,
% at the end of one line and again at the head of the next. As printed.
i det pragtfulde Guds Huus, naar Præsten prædiker ---
til Beroligelse!“ . . . . . „For Alvor vil man Intet höre
om det Forfærdelige, man vil lege det efter, omtrent,
som naar i Fredstider Krigerne eller rettere Ikke-Krigerne
lege Krig; man vil kunstnerisk fordre Alt af Taleren,
men selv vil man verdsligt og ugudeligt i Guds Huus
sidde ganske tryg, fordi man jo nok veed, at den Magt
har ingen Taler, den Magt, som kun Styrelsen har, at
gribe et Menneske, at kaste ham hen i Omstændighe\-%
dernes Vold, at lade Tilskikkelser og Prövelser og An\-%
fægtelser i Alvor prædike for ham til Opvækkelse.“

Efter de Forestillinger, man fra flere Sider har for\-%
sögt at udbrede iblandt Folket, skulde Udenforstaaende
nu vente, at Talens Öiemed altsaa var, at opskrue
Prædikeforedragets Strenghed til en saadan Höide, at
Folk derved nödvendigviis maatte afskrækkes fra at
söge Kirken. Enhver, der har læst Talen, veed, hvor
lös, hvor overfladisk og falsk denne Forestilling er. Vel
er Talen streng, men dog kun streng, forsaavidt den
religiöse Fordring i sig selv er streng; her er i Talen
ingen Overdrivelse; her, som overalt i det Religiöse, er
den uendelige Strenghed igjen væsentlig Mildhed.
% --- p. 14 ---
% ö/ø: Besöget, ifölge, gjöre, Selvprövelse, Indövelse, henföres, förste,
% förer — all o-DIAERESIS. No ø.
% PRINTER'S DEFECT: the first em-dash below (after „Troen“) is a battered
% sort, part of the rule not inked, so that it prints like a short broken
% double stroke. Set as an em-dash.
„Naar Troen“ --- saaledes slutter Talen --- „sees fra
dens ene, den himmelske Side, da seer man kun Salig\-%
hedens Gjenskin i den; men seet fra dens anden, den
blot menneskelige Side, seer man idel Frygt og Bæven.
Men saa er jo ogsaa den Tale usand, der idelig og
aldrig anderledes end indbydende, lokkende, vindende
vil tale om Besöget i Herrens Huus, thi fra den anden
Side seet, er det forfærdeligt. Derfor er ogsaa den Tale
usand, der tilsidst endte med ganske at skrække Men\-%
neskene bort fra at komme i Herrens Huus, thi fra den
anden Side seet er det saligt, een Dag i Guds Huus
bedre end tusinde ellers“. . . . . „Thi til Taleren hedder
det: brug al den Dig forundte Evne, villig til enhver
Opoffrelse og Eftergivenhed i Selvfornegtelse, brug den
til at vinde Menneskene --- men vee Dig, om Du vinder
dem saaledes, at Du udelader Forfærdelsen; brug derfor
al den Dig forundte Evne, villig til ethvert Offer i Selv\-%
fornegtelse, brug den til at forfærde Menneskene --- men
vee Dig, om Du ikke bruger den i Grunden dog, for at
vinde dem for Sandheden.“

Da Christendommen, ifölge Kierkegaards Opfattelse,
ikke er, hvad man saadan uden videre kalder en Lære,
men i væsentlig Forstand „en Existensmeddelelse“, var
det ham især magtpaaliggende at gjöre det saa tydeligt,
som muligt, hvorledes den Talendes egen Existens, uden
at Fordringens Strenghed gav det Menneskelige Paaskud
til at unddrage sig Forpligtelsen, og endda frisk væk be\-%
raabe sig paa „Troen“ og paa „den rene Lære“, kunde
og burde forholde sig til „Existensmeddelelsen“. I denne
Henseende faaer endog den forskjellige Maade, hvorpaa
han lod sine Skrifter udgaae, en mere end tilfældig Be\-%
% LETTERSPACED: the two book-titles „Til Selvprövelse“ and
% „Indövelse i Christendom“.
tydning. Sit Skrift: „\emph{Til Selvprövelse}“ udgav han
f.\ Ex.\ i sit eget Navn, „\emph{Indövelse i Christen\-%
dom}“, henföres derimod, i förste Udgave, til en Pseudo\-%
nym. Dette er nu ret betegnende for hvad han kalder
sin „Taktik“. Til det Sprog „Anti-Climacus“ förer i
% --- p. 15 ---
% ö/ø: Indövelsen, hidrörende, Höieste, höres, bör (twice), gjöre,
% Indrömmelse, Först, Selvprövelse, nöiagtig, höiærværdig, överste,
% tilhörer, Inderliggjörelse, utilböielig, rörer, förste — all o-DIAERESIS.
% No ø.
„Indövelsen“ svarer en saa anstrengt Existens, at For\-%
fatteren just af den Grund har stillet sig bag Pseudo\-%
nymen, og kun anmeldt sig som Udgiver: „I dette
Skrift, hidrörende fra Aaret 1848, er Fordringen til det
at være Christen tvunget op til et Idealitetens Höieste.
Dog siges, fremstilles, höres bör jo Fordringen, der bör,
christeligt, ikke slaaes af paa Fordringen, ei heller den
forties --- istedetfor at gjöre Indrömmelse og Tilstaaelse
sig selv betræffende.“ Det Skrift derimod, han ligefrem
har paatrykt sit Forfatternavn, skal ved den Aand og
Tone, hvori Fremstillingen fra Först til Sidst er holdt,
minde Læseren om det Afstandsforhold, hvori han selv
befinder sig til Idealerne.

„Lad mig“, siger han (Til Selvprövelse, S.\ 16),
„nöiagtig bestemme, hvor jeg, saa at sige, er. Der er
iblandt os en höiærværdig Olding, denne Kirkes överste
Geistlige; det han, hans Prædiken har villet, det samme
er det, jeg vil, kun en Tone stærkere, hvad der ligger
i min Personligheds Forskjellighed, og hvad Tidens For\-%
skjellighed kræver“. . . . . „Væsentligen tilhörer jeg Gjen\-%
nemsnittet, og her er det jeg har arbeidet for Uro i
Retning af Inderliggjörelse“.

„Thi christelig er der to Arter sand Uro. Uroen i
Troesheltene og Sandhedsvidnerne, hvilken sigter til at
reformere et Bestaaende. Saa langt har jeg aldrig vovet
ud; det er ikke Noget for mig, og forsaavidt Nogen i
Samtiden kunde synes at ville vove saa langt ud, var
jeg ikke utilböielig til at tage Polemik imod ham, for at
bidrage til, at det maatte blive aabenbart, om han er
den Berettigede“. . . . .

„Istedetfor selv tomt og opblæst at udgive mig
for et Sandheds-Vidne og foranledige Andre til dumdri\-%
stigt at ville være det Samme, er jeg en umyndig Digter,
der rörer ved Hjælp af Idealerne.“

Saavidt om „Forfattervirksomheden“ i det förste Af\-%
snit. Lykkes det nu den uhildede Iagttager ved disse
% --- p. 16 ---
og lignende Træk at danne sig en nogenlunde „sikker
Dom“, om S.\ Kierkegaards Grundtanker og Bestræbelser
i dette förste Afsnit af hans offentlige Liv, da vil en
„sikker Dom“ om hans „mentale Tilstand“ i det sidste
Afsnit falde som af sig selv. Som Forudsætningerne ere,
ere ogsaa Conseqvenserne, at sige, naar man har en Mand
for sig, som vil, hvad det end koster, være conseqvent.
At indrömme Forudsætningerne er altsaa med det Samme
at vedkjende sig Conseqvenserne, at sige, hvis Indröm\-%
melsen beroer paa Indsigt, ikke paa en flygtig Stemning.

Dersom det nu altid var en saa ganske ligefrem
Sag, at Conseqvensernes Overeensstemmelse med deres
Forudsætninger strax maatte falde endog en overfladisk
% „Öinene“: ö (diaeresis) confirmed at 600 ppi. No ø anywhere on p. 16.
Betragter i Öinene, vilde udförligere Oplysninger her
være overflödige. Men Enhver, der kjender Noget til
Livets Conflicter, veed med sig selv, at det Tilfælde
endda just ikke saa sjældent indtræffer, da det faaer
Udseende af conseqvent Holdning, hvad dog i Grunden
er jammerlig Inconseqvens, medens omvendt, det som
efter Skinnet ligner et Brud paa al Conseqvens, er i sit
Væsen det eneste conseqvente. Jo mere Inderlighed
Sagen har, desto smerteligere ville Collisionerne föles;
er det Inderlighedsproblemerne selv i deres Maximum,
hvorom det gjælder, vil Knuden snart strammes i den
Grad, at det kan blive Alvor med det berömte Sted hos
% German quoted in the same antiqua; „lügen“ with a true ü.
Jacobi: „Ja, ich bin der Atheist und Gottlose, der, dem
Willen, der Nichts will, zuwider, lügen will, wie Desde\-%
% FIVE printed points (the first close up to „log“, then four spaced).
mona sterbend log. . . . .“

Som Bindeled imellem förste og sidste Afsnit af vor
Forfs.\ offentlige Liv have vi „Indövelse i Christendom“. Til
det nye Oplag af dette Skrift har K.\ (Fædrel.\ 16.\ Mai 1855
Nr.\ 112) sendt et veiledende Ord. „Denne Bog har jeg“,
siger han, „ladet udgaae i et aldeles uforandret Oplag, fordi
% FIVE points here and five again at „Maade“ below.
jeg betragter den som et historisk Actstykke. . . . . Min
tidligere Tanke var: skal det Bestaaende forsvares, er
dette den eneste Maade. . . . . Paa den Maade kommer
% --- p. 17 ---
% Sheet signature „2“ at the foot of this page: not transcribed.
der dog Sandhed ind i det Bestaaende, det forsvarer
sig ved at dömme sig selv, det vedgaaer den christelige
Fordring, gjör Tilstaaelse sig selv betræffende om sin
% FIVE points.
Afstand. . . . . Dette var i mine Tanker den eneste Udvei
til, christeligt, at forsvare det Bestaaende, og for ikke
at forskylde paa nogensomhelst Maade at gaae for hur\-%
tigt tilværks, vovede jeg at give Sagen den Vending, for
at see, hvad den gamle Biskop vilde gjöre derved.
Hvis der var Kraft i ham, maatte han gjöre Eet af To:
enten afgjort erklære sig for den Bog, vove at gaae med
den, lade den være Forsvaret, der afværger, hvad Bogen
digterisk indeholder, Sigtelsen mod hele den officielle
% Nested quotation reusing the same „…“ pair, as everywhere in this article.
Christendom, at den er en Öienforblændelse „ikke en
suur Sild værd“; eller saa afgjort som muligt kaste sig
mod den, stemple den som et gudsbespotteligt og van\-%
helligt Forsög, og erklære den officielle Christendom for
den sande Christendom. Han gjorde ingen af Delene,
han gjorde Intet; han saarede blot sig selv paa Bogen,
og mig blev det tydeligt, at han var afmægtig.“

Det er ikke her Stedet at imödegaae alle Indven\-%
dinger af alle Slags, som ere reiste mod S.\ Kierkegaard
paa det Sidste. Skjöndt Beskyldningerne ere mange, er
Skylden dog kun een: Conseqvensen! For Conseqven\-%
sens Skyld ville vi her indskrænke os til to fremragende
Punkter: „Angrebet paa Biskop Mynster“, og Angrebet
paa det Bestaaende.

% The footnote is marked *) in the printing, set smaller under a short
% left-hand rule at the foot of the page; the mark stands after „1854“.
Da Dr.\ Sören Kierkegaard, Mandagen den 18de De\-%
cember 1854 ved sin Artikel\footnote{Artiklen skreven i Februar 1854.} i „Fædrelandets“ Nr.\ 295,
% LETTERSPACED DISPLAY. The Sperrsatz begins at „Var“ and runs, unbroken
% and through both quotation pairs and the dash, to the question mark;
% „vare der vistnok“ resumes ordinary spacing on the same line.
begyndte med det Spörgsmaal: \emph{Var Biskop Mynster
et „Sandhedsvidne“, et af „de rette Sandheds\-%
vidner“, --- er dette Sandhed?} vare der vistnok
Mange, der forargedes, Mange, der raabte paa Inconse\-%
qvens, Mange, der tænkte paa „herostratisk Berömmelse
ved at udslynge Fornærmelser“ o.\ s.\ v., men, som det
% --- p. 18 ---
lod til, ikke ret Mange, der besindede sig paa, at den
Opsigt vækkende Indsigelse i det Væsentlige, hvad den
afdöde Biskop angaaer, ikke indeholdt noget Nyt, ikke
indeholdt Andet, end hvad K.\ allerede forud, om end i
mildere Tone, havde sagt Biskop Mynster selv i hans
% „Öine“ here is ö, not ø — checked at 600 ppi against the ø of p. 19.
aabne Öine offentlig og saa tydelig, at Vedkommende
ikke kunde misforstaae det, og det uden at enten Bi\-%
skoppen selv eller Nogen iblandt hans Tilhængere den\-%
gang saameget som gjorde Mine til at have Noget der\-%
imod at indvende.

Den offentlige og bestemte Indsigelse imod, at Biskop
Mynster anvises Plads blandt „de rette Sandhedsvidner“,
% LETTERSPACED: the title only, inside its quotation marks.
er at læse i Skriftet: „\emph{Til Selvprövelse}“ S.\ 16 og 17.
Efterat K.\ nemlig S.\ 16 med de allerede citerede Ord:
„Der er iblandt os en höiærværdig Olding, denne Kirkes
överste Geistlige; det han, hans Prædiken har villet,
det Samme er det, jeg vil, kun en Tone stærkere, hvad
% „min Persons Forskjellighed“ — so printed here, where the same sentence
% was quoted on p. 15 as „min Personligheds Forskjellighed“, and is quoted
% again with „Personligheds“ eleven lines below. Transcribed as printed.
der ligger i min Persons Forskjellighed“, har i Udtryk
saa betegnende, at Ingen kan tage feil af, hvem „den
höiærværdige Olding, denne Kirkes överste Geistlige“,
% No dash at the end of this line or the next: OCR invents one; the image
% shows the line ending flush after „Mynsters,“.
er, sammenlignet sin Virksomhed med Biskop Mynsters,
og altsaa ligefrem stillet sig ved hans Side, ja, med den
Tilföining: „en Tone stærkere, hvad der ligger i min
Personligheds Forskjellighed“, endog --- thi i slige Sager
gjælder det hverken om Höflighed eller Fornærmelse,
men om Sandhed --- om man saa vil, ved hans höire
Side; er det en Selvfölge, at ethvert Ord, hvormed K.
charakteriserer sin Afstand fra „Troesheltene“ og „Sand\-%
hedsvidnerne“, nödvendigviis ogsaa kommer til at gjælde
for Biskoppen. Naar han da siger om sig selv, at han
„væsentligen tilhörer Gjennemsnittet“, siger han det
Samme om Biskop Mynster; naar han om „Uroen i
Troesheltene og Sandhedsvidnerne, hvilken sigter til at
reformere et Bestaaende,“ siger: „Saa langt har jeg al\-%
drig vovet ud, det er ikke Noget for mig, og forsaavidt
Nogen i Samtiden kunde synes at ville vove saa langt
% --- p. 19 ---
% Sheet signature „2*“ at the foot of this page: not transcribed.
% THIS PAGE MIXES THE TWO SORTS: ø in „Følgeblade“ (l. 19) and „Øine“
% (l. 31), ö everywhere else, including „Öine“ eight lines above the first
% ø. Each verified at 600 and 1200 ppi.
ud, var jeg ikke utilböielig til at tage Polemik imod
% FOUR points here (not five).
ham“. . . .; da maa Enhver, der har lært at læse, uvil\-%
kaarlig forstaae det Sagte, som om der stod: „Saa langt
have vi, Biskop Mynster og jeg, aldrig vovet ud, det
er ikke Noget for os; og forsaavidt Nogen i Samtiden
kunde synes at vove saa langt ud, vare vi, Biskop
Mynster og jeg, ikke utilböielige til at tage Polemik
imod ham.“ Kort, hvergang Kierkegaard ved sine „Ind\-%
römmelser“ og „Tilstaaelser“ bukker for „Troesheltene“,
og böier sig for „Sandhedsvidnerne“, seer det i Læserens
Öine livagtig ud, som om han just holdt Biskop Mynster
ved Haanden, for derved at formaae denne Kirkens
„överste Geistlige“ til at gjöre det Samme.

Havde nu enten Biskoppen selv eller Nogen af hans
Nærmeste fölt sig overbeviist om, at her ved denne
Sammenligning var skeet et Overgreb; hvorfor blev det
da ikke sagt K.\ strax og lige ud? Man kunde jo, thi
% ø: „Følgeblade“ — the slash is unmistakable at 600 ppi.
som K.\ senere har sagt (Følgeblade til: Dette skal siges)
„ved saadanne Forretninger bruger man Ordene sandt
som en beskrivende Naturforsker, Complimenter have
ikke hjemme her,“ have sagt ham uden Omsvöb: „Ved
saaledes at sammenligne Dig med „denne Kirkes överste
Geistlige,“ er det ikke Biskop Mynster Du fornærmer,
men Du krænker Sandheden, idet Du ved en saa upassende
Sammenligning forvirrer Begreberne. Du vil klare Exi\-%
stenserne; hvad er vel din Existens? Du kan gjerne
være Forfatteren med „den rige Dannelse,“ „den dybe
religiöse Grund,“ „det Sind, som er mægtigt til at træde
% ø: „Øine“ — and „höieste“ on the very same line is ö.
Livets höieste Problemer under Øine,“ men din Existens?
Ja, deri har Du Ret, den „tilhörer Gjennemsnittet.“ Men
Biskop Mynsters Existens, saaledes mærket som den er,
med Byrdernes, Korsenes, Lidelsernes Tegn, tilhörer ikke
Gjennemsnittet; eller kan Du i selvsyg Forblindelse maa\-%
skee ikke see, at Biskop Mynster, foruden hvad han er i den
humane Retning, tillige er et Andet og Höiere, end de
Mænd, vi ellers regne blandt Landets Ypperste, at han
% --- p. 20 ---
just er et Sandhedsvidne, et af de rette Sandhedsvidner,
et Led i den hellige Kjæde?“

Havde man fra först af været saa aabenhjertig, saa
var K.\ naturligviis strax bleven orienteret; nu derimod,
da man iagttog Taushed, drömte han vist ikke om, at
det maaskee kun var en diplomatisk Taushed, men tog,
som man seer, Tausheden ligefrem for en stiltiende Ind\-%
römmelse. „Saaledes var,“ siger han i Indsigelsen, „naar
man lægger det nye Testamente ved Siden, Biskop Mynsters
Christendoms-Forkyndelse en, især for et Sandhedsvidne,
mislig Christendoms-Forkyndelse. Men da var der, mente
jeg, det Sande hos ham, at han, hvad jeg holder mig
fuldt forvisset om, var villig til at tilstaae for Gud og
sig selv, at han ingenlunde, ingenlunde var et Sandheds\-%
vidne, --- just denne Tilstaaelse var i mine Tanker det
Sande.“

Hvor er altsaa det Inconseqvente? Havde K.\ tidligere
begyndt med at complimentere Biskop Mynster som „Sand\-%
hedsvidne,“ fordi han nu af visse Grunde saa godt kunde
lide Biskop Mynster; men sidenefter reist en „herostratisk“
Indsigelse imod at kalde Biskop Mynster et Sandheds\-%
vidne, fordi han nu fölte sig fornærmet, og ikke længere
kunde lide Biskop Mynster: det havde været inconse\-%
qvent. Havde han derhos, for at dække sin Inconseqvens,
gjort Qvakleri med Ordet: Sandhedsvidne, at det jo f.\ Ex.\ baade
kunde betyde Dit og betyde Dat, at han vistnok
tidligere var uforvarende kommen til at bruge det Udtryk:
Sandhedsvidne, men dog paa ingen Maade havde meent,
at man derved skulde tænke sig: det rette Sandhedsvidne,
nu derimod, da han just vilde paapege, hvad man skulde
tænke sig ved det rette Sandhedsvidne, faldt Biskop
Mynster igjennem: da havde en saadan Adfærd ikke
blot været inconseqvent, men den havde været foragtelig.
Nu derimod, da det for Gud og hver Mand er vitterligt,
at Kierkegaard selv i sidste saavelsom i förste Afsnit af
sit Liv fastholder det samme, det selvsamme Begreb om
% --- p. 21 ---
„Sandhedsvidnet,“ og i Forhold til det at være Sandheds\-%
vidne giver os i Grundtræk det samme, det selvsamme
Billede af Biskop Mynster: hvori bestaaer saa det In\-%
conseqvente? hvad er det, som i denne hans Fremgangs\-%
maade skal egne sig til at vække Mistanke mod hans
„mentale Tilstand“?

% NOT letterspaced here, though the same title is letterspaced on p. 18
% and again, unspaced, on p. 27. Checked on the image.
Hvorledes Kierkegaard vilde have det i sit anförte
Skrift (Til Selvprövelse S.\ 17) givne Tilsagn, at han, for\-%
saavidt Nogen i Samtiden kunde synes at ville vove „saa
langt ud“ (at han nemlig i Lighed med „Sandheds\-%
vidnerne“ vilde reformere et Bestaaende) ikke var util\-%
böielig til „at tage Polemik imod ham,“ forstaaet, har han
viist i sin Artikel (Fædrel.\ 31.\ Jan.\ 1851 No.\ 26) mod
Dr.\ Rudelbach. Artiklen begynder saaledes: „Yttringen fin\-%
des“ (i Dr.\ Rudelbachs Skrift om det borgerlige Ægteskab
1851) „p.\ 70. I Texten staaer: Sandelig, Kirkens dybeste
% FOUR points.
og höieste Interesse i vore Dage er. . . . netop at blive
% Suspended hyphen kept as printed.
emanciperet fra det, man med Rette kan kalde Vane- og
Statschristendom. Hertil saa Noten 121. Det er det
samme som en af vore meest udmærkede Skribenter fra
de sidste Dage, Sören Kierkegaard, söger at indprente,
indpræge, og, som Luther siger: inddrive alle dem, som
ville höre.“

Denne Artikel er i Overeensstemmelse med Kierke\-%
gaards Princip affattet i streng conservativ Aand. „Intet,“
hedder det, „er jeg saa betænkelig ved som ved Alt,
hvad der endog blot smager af denne usalige Forvexling
af Politik og Christendom, en Forvexling, der saa let
kan bringe en ny Art Kirkereformation op og i Mode,
% Both OCR witnesses read a dash before „den“; at 600 ppi the mark is a
% speck in the left margin, well outside the measure and a fraction of a
% dash's width. Not transcribed.
den bagvendte Reformation, der reformerende sætter et
nyt Slettere i Stedet for et gammelt Bedre, medens det
dog skal være vist og sandt, at det er en Reformation,
i hvilken Anledning hele Byen bliver illumineret. Christen\-%
% „Inderligjörelse“ with one g, as printed; „Inderliggjörelse“ with two on
% p. 22. Both transcribed as they stand.
dom er Inderlighed, Inderligjörelse. Ere til en given
Tid de Former, hvorunder der skal leves, ikke de fuld\-%
komneste: nu, i Guds Navn; kan man faae dem bedre:
% --- p. 22 ---
% FOUR points.
vel. Men væsentligen: Christendom er Inderlighed.“ . . . .
„Skulde dette höre med til sand Christendom, denne Tro
paa de politisk opnaaede, frie Institutioners frelsende
Magt, saa er jeg ingen Christen, ja endnu værre, jeg er
et reent Djævelens Barn; thi, oprigtig talt, jeg har end\-%
og Mistanke til disse, politisk opnaaede, frie Institutioner,
især til deres frelsende og gjenfödende Magt. Det er
% PRINTER'S ERROR: „christendum“ set as one word, the word-space dropped
% (Kierkegaard's text reads „saa christen dum er jeg“). As printed.
min Christendom, eller saa christendum er jeg, som da
forresten aldrig har befattet mig med „Kirke“ og „Stat“
--- det er meget for stort for mig . . . . Jeg har heller
aldrig kjæmpet for „Kirkens“ Emancipation, saa lidet som
for den grönlandske Handels, Qvindernes eller nogen\-%
somhelst anden Emancipation. Jeg har som Enkelt, med
Sigtet „den Enkelte,“ tilsigtende Inderliggjörelse i Christen\-%
dom i „den Enkelte,“ conseqvent, med Aandens Vaaben,
ene og alene med Aandens Vaaben, kjæmpet for at gjöre
opmærksom paa og mod at lade sig bedrage af „Sandse\-%
bedrag“ . . . . „Det er ikke udvortes fra Christendommen
skal hjælpes, ved Institutioner og Constitutioner, og da
allermindst, naar disse ikke, paa gammelt Christeligt,
skulde lidende tilkjæmpes ved Martyrier, men paa Politisk,
% Printed with a thin space on each side of the hyphen
% („selskabeligt - venskabeligt“) in this stretched line; closed up here as
% justification, per the convention for spaced punctuation.
selskabeligt-venskabeligt vindes ved Ballotation eller i
Tallotteriet --- at hjælpes paa den Maade er tvertimod
dens Undergang.“

Dersom Nogen imidlertid i disse Yttringer af et „con\-%
servativt“ Sindelag mener at finde et Beviis for, at
Kierkegaard herved saa ligefrem havde forskrevet sig
til det Bestaaende, at han ikke, uden udtrykkelig at bryde
med sine egne Principer, nogensinde kunde tillade sig
at angribe det Bestaaende; da vil en Saadan, ved at
see lidt nærmere til, let overbevise sig om, at han har
taget mærkelig feil.

„Lad mig dog,“ hedder det som i en Efterskrift til
Artiklen, „tilföie Fölgende, for at, hvad jeg siger, ikke
skal blive misforstaaet, som var det min Mening, at
Christendom ene og alene bestod i at finde sig i Alt be\-%
% --- p. 23 ---
% ø on this page: „Øiet“ only; every other o-diacritic is ö.
træffende de ydre Former“ . . . . „I Apostlernes Gjer\-%
ninger læse vi de Ord: man bör adlyde Gud mere end
Mennesker. Altsaa der gives Tilfælde, hvor et Bestaaende
kan være af den Beskaffenhed, at den Christne ikke bör
finde sig deri, ikke bör sige, at Christendom just er
denne Ligegyldighed for det Udvortes.“

„Men lad os nu see, hvorledes Apostlerne ikke bare
sig ad, --- thi hvorledes disse ærværdige Skikkelser bare
sig ad, det veed vel Enhver.“ „Apostlerne gik ikke saa\-%
dan og snakkede med hinanden og sagde: „Det er
utaaleligt, at Synedriet sætter Straf paa Ordets For\-%
kyndelse, det er Samvittighedstvang. Dog, hvad skulle
vi gjöre? Skulle vi ikke see at blive nogle Stykker, og
saa indgive en Adresse til Synedriet, eller see at komme
med paa en Synode; det var ikke umuligt, at vi saa
ved at holde sammen med hvad der ellers er vore Fjender,
kunne ved Ballotation faae Majoriteten, og vi da kunne
faae Samvittighedsfrihed til at forkynde Ordet.“ Gud i
Himlene! Ærværdige Skikkelser, tilgiver, at jeg har maat\-%
tet tale saaledes, det var nödvendigt.“

„Hvorledes bare de sig derimod ad --- thi der er
dog vist Adskillige, der have glemt det.“ . . . . „Apost\-%
len er væsentligen en eenlig Mand.“ . . . . „Lad os an\-%
tage, at En möder ham, der siger: Veed Du, at Syne\-%
driet har sat Hudflettelse paa at forkynde Ordet. Apost\-%
len svarer: Naa, har Synedriet det, saa vil jeg altsaa
blive hudflettet. Imorgen sætter Synedriet Dödsstraf.
Apostlen svarer: Naa, har Synedriet gjort det, saa vil
jeg altsaa blive henrettet. Han lader det Bestaaende
bestaae; ikke et Ord, ikke en Stavelse, ikke et Bogstav
i Retning af Forandring i det Udvortes, ikke den flygtigste
% ø: „Øiet“.
Tanke i hans Hoved, ikke et Blink med Øiet, ikke en
% „Apostelen“ here, „Apostlen“ four times above. As printed.
Mines Bevægelse i denne Retning. „Nei,“ siger Apostelen,
„lad det kun staae urokket fast, dette Bestaaende, thi
dette staaer med Guds Hjælp ogsaa urokket fast, at idag
bliver jeg hudflettet, imorgen henrettet, eller, hvad der
% --- p. 24 ---
er det Samme, idag forkynder jeg Ordet, og imorgen,
Amen.“

„Og her et christeligt Memento. Ordenligviis er
Christendom indenfor Christenhed just Ligegyldighed for
det Udvortes i Selvbekymring. Men colliderer En saa\-%
ledes med det Bestaaende, at det kunde falde ham ind,
% „store Gud“ with a plain o here; the same phrase is reset on p. 28 with
% an ø sort (see the note there).
at det er et Samvittighedsspörgsmaal, --- store Gud, et
Samvittighedsspörgsmaal! --- og han vover at sige dette:
da har han at være eenlig Mand, at stride lidende, at
vælge Martyriet. Samvittighed og Samvittighedssag kan
--- dette er evig vist og ligger i Sagen selv, thi Sam\-%
vittighed er ingen numerisk Bestemmelse --- kun re\-%
præsenteres af eenlig Mand og i Charakteer, ved Hand\-%
ling.“ . . . .

Da denne Artikel --- just en Artikel efter Biskop Myn\-%
sters Hjerte --- saae Lyset, var der neppe Nogen iblandt
dem, der dengang holdt paa det Bestaaende, som heri
fandt noget Stödende, end sige da, noget saa Stödende,
at det i fjerneste Maade kunde vække Mistanke mod
Forfatterens „mentale Tilstand.“ De i Artiklen udtalte
Grundsætninger ere da ogsaa virkelig, det maa Enhver,
der vil bruge sin Forstand, indrömme, af den Natur, at
baade Staten og Christendommen, at sige, det nye Te\-%
stamentes Christendom kunne, hver efter sin Interesse,
vedkjende sig dem. Det, der i religiös Henseende for\-%
aarsager Staten de mange Plagerier, Sammenstöd og uhyg\-%
gelige Bryderier, er ingenlunde „eenlig Mand, der væl\-%
ger at stride lidende“, men meget mere det politiserende
Kirkevæsen, den religiöse Verdslighed, der dog vil over
det Verdslige, den Samvittighedernes Selskabsdrift, der
grunder sin Frihed paa „Ballotation.“ „Alt, hvad der
er Parti og vil virke som Parti, maaskee endog intrigue\-%
rende --- naar det mod et Bestaaende vil beraabe sig
paa „Samvittighed“, da forskylder det en Usandhed“, er
en Grundsætning, som, mærkværdigt nok, i lige Grad, om
end i modsat Aand, kan hyldes af Statsmænd og Apostler.
% --- p. 25 ---
Spörge vi nu, hvorvidt Kierkegaards, som det synes,
i höieste Grad agiterende, voldsomme og revolutionære
Optræden i de sidste Afsnit af hans Liv, uden Omsvöb,
undskyldende Vendinger eller spidsfindige Kunstgreb la\-%
der sig bringe i Overeensstemmelse, i grundig, sand og
fuldstændig Overeensstemmelse med Tankegangen i nys\-%
nævnte fredsommelige Artikel, er Svaret ikke langt borte.
Kun maae vi, hvad de sidste Optrin angaaer, ikke glem\-%
me, at Buen er spændt, ikke glemme, hvad det var for
et höittravende „Vidnesbyrd“, der gav Anledning til, at
Buen saaledes blev spændt, ikke glemme, at Inderlig\-%
hedsknuden er strammet, at det nu er Alvor, virkelig
% The German is NOT enclosed in quotation marks here.
Alvor med „det berömte Sted“ hos Jacobi: Ja, ich bin
der Atheist und Gottlose!

% LETTERSPACED DISPLAY: „Exempel:“ is in ordinary spacing; the Sperrsatz
% opens at the quotation mark and closes after „Anden.“, the bracketed
% source that follows being ordinary again. The quotation is never closed
% — there is no “ before the parenthesis — so it stands +1 in the balance.
% The second footnote of the article, mark *), three lines under a short
% left-hand rule, its last four words letterspaced.
Exempel: \emph{„Hvad der skal gjöres --- det skee
nu ved mig eller ved en Anden.} (Anden Pintse\-%
dag 1854.\footnote{Her kan gjentages, hvad K.\ selv udtrykkelig fremhæver i en
tidligere Artikel (Ved Biskop Mynsters Död): „Man bemærke
\emph{Aarstal og Datum}“.} S.\ Kierkegaard. Fædrelandet, 22.\ Marts 1855.
Nr.\ 69.)

„Der skal först og fremmest, efter den störst mulige
Maalestok, gjöres en Ende paa hele den officielle --- vel\-%
meente --- Usandhed, der --- velmeent --- fremkogler og
opretholder det Skin, at det er Christendom, man forkyn\-%
der, det nye Testamentes Christendom. I denne Hen\-%
seende gjælder det om ikke at skaane.“ . . . .

„Naar saa det er gjort, skal Sagen vendes saaledes:
er det sande Sammenhæng dog ikke egentlig dette, at
det, fra Slægt til Slægt, er gaaet saaledes til Agters for
os Mennesker, . . . . at vi, istedetfor den uforskammede
Sludder om, at Christendommen er perfectibel, at vi
gaae fremad, ja, at Christendommen maaskee ikke en\-%
% Suspended hyphen („Dusin- og Snese-Mennesker“) kept as printed.
gang mere tilfredsstiller --- at vi, usle, jammerlige Dusin- og
Snese-Mennesker, til syvende og sidst virkeligen ikke
% --- p. 26 ---
kunne bære dette Guddommelige, som er det nye Te\-%
stamentes Christendom; og at vi for saavidt maae nöies
med den Art Religiösitet, som nu er den officielle, efterat
det, vel at mærke, er gjort kjendeligt, at den ingenlunde
% FOUR points here, THREE at „Salighed“ five lines below.
er det nye Testamentes Christendom. . . . Saaledes maa
Sagen vendes; bort, bort, bort med al Öienforblindelse,
frem med Sandheden, ud med Sproget: vi duer ikke til
at være Christne i det nye Testamentes Forstand, --- og
dog trænge vi til at turde haabe en evig Salighed.“ . . .
„Naar saa Sagen er vendt saaledes, saa vil det vise sig,
om der er noget Sandt i dette, om det har Styrelsens
% The dash here prints in two pieces, the middle of the sort having failed;
% it is an em dash, set as one.
Samtykke, --- hvis ikke, saa maae Alt briste, at der i
denne Rædsel atter maatte blive Individer til, som kunne
bære det nye Testamentes Christendom. Men Ende,
Ende skal der gjöres paa den officielle --- velmente ---
Usandhed.“

Men er det dog ikke oprörende? „Ja, ich bin der
Atheist und Gottlose, der dem Willen, der Nichts will, zu\-%
wider, lügen will.“ Havde K., istedetfor at indlade sig paa
saa „radicale“ Fordringer og „pessimistiske“ Anskuelser,
kunnet beqvemme sig til at slaae ind paa en jævnere Vei;
havde han, træt af sin Stilling som „eenlig Mand“, om\-%
sider fulgt den christelige Selskabsdrift, eller, for at vælge
et bedre Udtryk, været saa „kjærlig“ at være som Andre,
og virke som Andre: det vilde være gaaet ham ganske
anderledes. Havde han, istedetfor at ivre mod „den
officielle Christendom“, og vække Skandale, viist den
Besindighed at fremkomme f.\ Ex.\ med et Forslag til Ordning
% „det skulde sikkert engang være lykkedes ham“ — as printed.
af den forventede Kirkeforfatning; det skulde sikkert en\-%
gang være lykkedes ham „at komme med paa en Sy\-%
node“; og var han först kommen med paa Synoden, og
havde faaet Ordet paa Synoden, maaskee endog de fleste
Stemmer ved „Ballotationen“, da skulde vist ingen for\-%
sigtig Mand have vovet at sige höit om S.\ Kierkegaard,
at han nu handlede inconseqvent, eller tillade sig at ind\-%
vende Noget mod hans „mentale Tilstand.“
% --- p. 27 ---
Den jævne Vei var indbydende nok; hvorfor valgte
han da ikke den jævne Vei?

At blænde de Tankelöse er en fattig Kunst; men
Enhver, som vil see, maa kunne see, at K.\ ikke uden
at bryde med sine Forudsætninger kunde indlade sig paa
Reformer ad fredelig Vei. For ham var der kun Valg
mellem Eet af To: enten blive staaende, som hidtil, paa
den Sætning: „Ordentligviis er Christendom indenfor
„Christenhed“ just Ligegyldighed for det Udvortes i Selv\-%
bekymring,“ eller, naar Tingenes Gang tog ham dette
„Ordentligviis“ bort, da saa resolut, saa ubetinget, saa
afgjörende, som muligt, angribe det Bestaaende.

„Saa omhyggeligt som der hidtil er blevet skjult
over, hvilken min Opgave kunde blive, saa forsigtigt
som jeg er bleven bevaret i den uigjennemtrængeligste
Ukjendelighed: saa afgjörende skal jeg ogsaa nu, da
Öieblikket er kommen, gjöre det kjendeligt. Det Spörgs\-%
maal, hvad Christendom er, og dermed igjen Spörgs\-%
maalet om Stats-Kirke, hvad man nu vil kalde det, Sam\-%
mensætningen eller Sammenholdet af Kirke og Stat skal
bringes til den yderste Afgjörelse. Det kan og skal ikke
% THREE points.
gaae, som det Aar efter Aar gik.“ . . .

Men er det da ikke aabenbart Oprör? Naar K.\ i
% NOT letterspaced here (cf. p. 18, where the same title is).
Skriftet: „Til Selvprövelse“, S.\ 16, spörger: „Saa forkynder
jeg da maaskee Tumult, Alts Omvæltning, Uorden?“ og
derpaa svarer: „Sandeligen nei“; kunde man da ikke
med Hensyn paa „Öieblikkene“ vende hans Ord som
Vaaben mod ham, og sige: „Sandeligen jo. Nu forkyn\-%
der du Tumult, Alts Omvæltning, Uorden!“? Denne Sag
er ikke indviklet. Det kierkegaardske Sprog er, uden
Omsvöb, revolutionært. „Thi ikke sandt, Du menige
% This dash prints broken in the middle — the sort has lost its waist;
% an em dash, set as one.
Mand, det kan Du meget godt forstaae --- jeg mener, at
just Du kan langt lettere og bedre forstaae det, end de\-%
% The dash at the end of this line prints with a vertical shoulder rising
% from its right end (a damaged or foul sort); an em dash by its place and
% length, set as one.
moraliserede Præster og en fordærvet Fornemhed ---
ikke sandt, dette kan Du ypperligt forstaae, at Eet er
at forfölges, mishandles, hudflettes, korsfæstes, hals\-%
% --- p. 28 ---
hugges o.\ D., et Andet, comfortabelt indrettet, med Fa\-%
milie, jævnt avancerende, at leve af at skildre, at en An\-%
den blev hudflettet o.\ s.\ v.“ Dette Sprog er revolutio\-%
nært, ingen Agitator skal saa let byde over.

Desuagtet er Kierkegaards Optræden det Revolutio\-%
nære fuldstændig modsat; vi have i denne Henseende
uimodsigelige Beviser. „Colliderer En saaledes med det
Bestaaende, at det kunde falde ham ind, at det er et
% ø AND PRINTER'S ERROR: „støre“, an ø sort where „store“ is meant — the
% same Kierkegaard sentence is set „store Gud“ on p. 24 with a plain o.
% The slash is unambiguous at 1200 ppi and the two settings were compared
% side by side. Transcribed as printed.
Samvittighedsspörgsmaal, --- støre Gud, et Samvittigheds\-%
spörgsmaal! --- og han vover at sige dette, da har han
at være eenlig Mand, at stride lidende, at vælge Marty\-%
riet.“ Disse Ord har Kierkegaard under sin sidste Op\-%
træden samvittighedsfuldt stræbt at sætte i Charakteer;
vi kunne nok fremhæve denne Kjendsgjerning, uden at
vi just derfor behöve at tale höitravende, og compli\-%
mentere den Afdöde som Martyr, eller som „Sandheds\-%
vidne.“ K.\ var og blev under hele Striden en „eenlig
Mand“: der var Ingen, der hjalp ham, Ingen, der kunde
i endelig Forstand komme ham til Hjælp. Selv om der
havde været, hvad der dog nok ikke var, En eller An\-%
den iblandt os, som, med samme Sikkerhed paa Katego\-%
rierne, tillige havde havt Resoluthed og Mod til at vove
ligesaa langt ud; det var dog ikke blevet til endeligt
Sammenhold imellem de Tvende: hver af de Stridende,
om de end i bedste Forstaaelse havde stridt i samme
Aand, for samme Sag, skulde dog have stridt for sig,
som eenlig Mand. Men at en saadan Fremgangsmaade
er det stik Modsatte af den revolutionære, seer dog vel
Enhver. Dersom man overseer denne Modsætning, hvor\-%
ledes vil man da forklare sig, at Kierkegaards revo\-%
lutionære Sprog kunde blive hört af „den menige Mand“,
uden at Nogen i Mængden rörte saameget som en Finger
til Angreb paa det Udvortes? Eller vidste Kierkegaards
Ord maaskee ikke at finde Veien til Lidenskaberne?
Mangler hans Sprog maaskee det Opflammende, der el\-%
lers saa let bevirker en Bruusning? Eller kunde „den
% --- p. 29 ---
menige Mand“ i sin Eenfoldighed dog alligevel ikke rig\-%
tig forstaae ham, naar han sagde, „at Eet er at forfölges,
mishandles, hudflettes, korsfæstes, halshugges o.\ dsl.,
et Andet, comfortabelt indrettet, med Familie, jævnt
avancerende, at leve af at skildre, at en Anden blev
hudflettet?“ Eller var det maaskee Præstekirkens apo\-%
stoliske Myndighed, der imponerede den menige Mand?

At den kierkegaardske Indsigelse, den stærkeste som i
nyere Tid er gjort mod et bestaaende Kirkevæsen, en Ind\-%
sigelse, der, idet den löd hele Landet over med Gjenklang
% PRINTER'S DEFECT: „först“ — the t has lost the left arm of its crossbar
% and prints like f (or a long ſ). „Först“ two sentences below, and
% „förste“ in the last line of the article, both have a sound t.
i tusinde og atter tusinde Sjæle, dog först og sidst löd
saaledes, at hver Mand begreb, at det ikke var en Om\-%
væltning i det Udvortes, men en Forandring i det Ind\-%
vortes, hvorom det gjaldt, derfor have vi ikke blot Fol\-%
kets Besindighed og sunde Sands --- den sunde Sands
kunde dog i Sligt desværre let bedaares af en Agitator,
--- men væsentlig og kategorisk Kierkegaards Position som
„eenlig Mand“, at takke.

Hvad her nu forövrigt er opnaaet ved den kierke\-%
gaardske Indsigelse, kan ikke oplyses i denne Sammen\-%
hæng. Opgaven var i nogle Træk at paavise Conseqven\-%
% PRINTER'S ERROR: a full point stands before the comma after
% „Forudsætninger“ — a clean round stop on the baseline, the same size as
% the comma beside it, checked at 1200 ppi. As printed.
serne af hans Forudsætninger., den indre Overeensstem\-%
melse af hans Ord og Handling fra Först til Sidst. De,
som beklage, at de ikke kunne have „nogen sikker Dom“
om dette Menneskes „mentale Tilstand i det sidste Afsnit
af hans Liv“, bör alvorlig overveie, om Grunden hertil
dog alligevel ikke skulde ligge hos dem selv, at de nem\-%
% PRINTER'S ERROR: „komme“ for „komne“. As printed.
lig i deres Forskninger endnu ikke ere komme saavidt,
at de kunne have nogen sikker Dom om „dette Menne\-%
skes“ Forfattervirksomhed i det förste Afsnit.

% End of the article. A short centred rule, then about a third of a page of
% white: no signature, no dateline, no „Fortsættes“.
\begin{center}\rule{0.25\linewidth}{0.4pt}\end{center}

\end{document}
